\documentclass[10pt,a4paper]{amsart}
\usepackage[margin=19mm]{geometry}
\usepackage{amsmath,amssymb,mathtools}
\usepackage[scr=rsfs,cal=euler]{mathalfa}
\usepackage{xfrac}
\usepackage[unicode,hidelinks,bookmarks=false]{hyperref}
\newcommand{\ie}{i.e.\ }
\newcommand{\cf}{cf.\ }
\newcommand{\resp}{respectively\ }
\newcommand{\Subsection}{\subsection}
\providecommand{\nobreakdash}{\nobreak\hbox{-}}
\setlength{\emergencystretch}{2em}
\multlinegap0pt
\usepackage{amssymb}
\usepackage{braket}
\newtheorem{lemma}{Lemma}
\newcommand{\R}{\mathbb{R}}
\renewcommand{\phi}{\varphi}
\title{Computing the free energy of quantum Coulomb gases and molecules via quantum Gibbs sampling}
\author{Simon Becker, Cambyse Rouz\'e, Robert Salzmann}
\date{Source: arXiv:2604.15263v1 (CC BY 4.0)}
\begin{document}
\maketitle

\noindent\emph{Source and licence.} Computing the free energy of quantum Coulomb gases and molecules via quantum Gibbs sampling. Version arXiv:2604.15263v1 (CC BY 4.0), distributed by arXiv under the Creative Commons Attribution 4.0 International licence, http://creativecommons.org/licenses/by/4.0/, as recorded in the arXiv licence metadata for this exact version and shown on the abstract page https://arxiv.org/abs/2604.15263v1. Full licence notice: \texttt{licenses/arxiv-2604.15263-CC-BY-4.0.txt}.\par\smallskip

\noindent\textit{Editorial scope.} Counted unit: the article's lemma lem:product-cutoff-riesz-no-scaling (source0.tex lines 1192--1200). The article's complete original proof of it is delivered with it (source0.tex lines 1202--1368, one \texttt{proof} environment of 167 lines). No mathematics is written, repaired or re-derived: the statement and the proof are the article's own lines, in the article's own notation, and no retained line is substituted or re-flowed.  No further source-local context is needed: every cross-reference of every retained line resolves inside the two delivered spans.  Nothing was omitted from any retained span, so the package carries no omission record.  No retained line contains a citation, so the wrapper carries no bibliography and no bibliography entry.  No retained line contains a figure environment, an image-inclusion command or drawing code, so no image is delivered and the package declares no asset dependency.  Editorial formatting only, in addition: an AMS \texttt{amsart} preamble that restates the article's own declarations and macros used by the retained lines, namely the environments \texttt{lemma} on separate counters, together with the standard packages those lines need.  The retained environments print in this document as Lemma 1 in order of appearance, whereas the article prints the same items with the numbering of its own full document.  The complete native source of the article as distributed in the arXiv e-print for this exact version is bundled unchanged as \texttt{sources/198594/source0.tex}.
\begin{lemma}[Product cutoff estimate without scaling assumption]
\label{lem:product-cutoff-riesz-no-scaling}
Given $d\in\{2,3\}$, for all \(n,M\) and \(\psi\in Q(H_{0,n})\),
\begin{align}
|W_n[\psi]|&\lesssim \,B_n\langle\psi,(H_{0,n}+1)\psi\rangle,\label{eq111}\\
|W_n[\Pi_M^\perp\psi]|&\lesssim\,B_n(\lambda_{M+1}+1)^{-1/2}\langle\psi,(H_{0,n}+1)\psi\rangle,\label{eq222}\\
|\Delta_{n,M}[\psi]|&\lesssim\,B_n(\lambda_{M+1}+1)^{-1/4}\langle\psi,(H_{0,n}+1)\psi\rangle.\label{eq333}
\end{align}
\end{lemma}

\begin{proof}
 Fix $1\le i<j\le n$ and $d=3$. Introduce the relative variable $r:=2^{-1/2}{(x_i-x_j)}$, so that ${|x_i-x_j|^{-1}}=2^{-1/2}{|r|^{-1}}$. By Hardy's inequality in $\mathbb R^3$,
\[
\left\|\frac{u}{|r|}\right\|^2\le 4\|\nabla_r u\|^2.
\]
Hence, by Cauchy--Schwarz and Young's inequalities, for every \(\eta>0\),
\begin{align*}
\left\langle u,\frac{1}{|x_i-x_j|}u\right\rangle
=
2^{-1/2}\left\langle u,\frac{1}{|r|}u\right\rangle
\le
\sqrt2\,\|u\|\,\|\nabla_r u\|
&\le
\eta\|\nabla_r u\|^2+ (2\eta)^{-1}\|u\|^2\\
&= \eta\|\nabla_r u\|^2+ C_\eta\|u\|^2,
\end{align*}
where we defined $C_\eta = (2\eta)^{-1}$ in the last equality.
Since
\[
\|\nabla_r u\|^2\le \langle u,(-\Delta_{x_i}-\Delta_{x_j})u\rangle
\le
\langle u,\bigl(-\Delta_{x_i}-\Delta_{x_j}+|x_i|^2+|x_j|^2\bigr)u\rangle,
\]
we obtain
\begin{equation}
\frac{1}{|x_i-x_j|}
\le
\eta\bigl(-\Delta_{x_i}-\Delta_{x_j}+|x_i|^2+|x_j|^2\bigr)+(2\eta)^{-1} \label{usefulestimate111}
\end{equation}
as quadratic forms. Taking \(\eta=1\), and using \(h_i+h_j+1\le 2(H_{0,n}+1)\), we get
\[
|W_n[\psi]|
\lesssim
\sum_{1\le i<j\le n}|\alpha_{n,i,j}|\langle\psi,(h_i+h_j+1)\psi\rangle
\lesssim
B_n\langle\psi,(H_{0,n}+1)\psi\rangle.
\]
In dimension \(d=2\), let \(u\in Q(h_1+h_2)\). Fix \(0<\varepsilon<1\). Since
\[
|\log r|\le C_\varepsilon\bigl(1+r^{-\varepsilon}+r^\varepsilon\bigr),\qquad r>0,
\]
it suffices to estimate these three terms separately. First,
\[
\int_{\R^2\times\R^2}|u(x,y)|^2\,dx\,dy=\|u\|^2
\le
\langle u,(h_1+h_2+1)u\rangle.
\]
Next, using the change of variables \(r=x-y\), \(s=x+y\), we obtain
\[
\int_{\R^2\times\R^2}|x-y|^{-\varepsilon}|u(x,y)|^2\,dx\,dy
=
\frac{1}{4}\int_{\R^2\times\R^2}|r|^{-\varepsilon}|u(r,s)|^2\,dr\,ds.
\]
Set
\[
F(r):=\|u(r,\cdot)\|_{L^2(\R^2_s)}.
\]
Then
\[
\int_{\R^2\times\R^2}|r|^{-\varepsilon}|u(r,s)|^2\,dr\,ds
=\int_{\R^2}|r|^{-\varepsilon}F(r)^2\,dr.
\]
Choose \(p>1\) such that \(\varepsilon p<2\), and let \(q=\frac p{p-1}\) its Hölder conjugate. Since
\(|r|^{-\varepsilon}\in L^p(B_1(\R^2))\) for $B_1(\mathbb{R}^2)$ the unit ball in $\mathbb{R}^2$, Hölder's inequality in the \(r\)-variable yields
\[
\int_{|r|\le1}|r|^{-\varepsilon}F(r)^2\,dr
\le
\bigl\||r|^{-\varepsilon}\bigr\|_{L^p(B_1(\mathbb{R}^2))}\|F^2\|_{L^q(\R^2)}
=
C\|F\|_{L^{2q}(\R^2)}^2.
\]
Now \(2q<\infty\), so by the Sobolev embedding
\[
H^1(\R^2_r;L^2(\R^2_s))\hookrightarrow L^{2q}(\R^2_r;L^2(\R^2_s))
\]
we obtain
\[
\|F\|_{L^{2q}(\R^2)}=\|u\|_{L^{2q}(\mathbb{R}_r^2,L^2(\mathbb{R}_s^2))}
\lesssim\|u\|_{H^1(\R^2_r;L^2(\R^2_s))}.
\]
Hence
\[
\int_{|r|\le1}\int_{\R^2}|r|^{-\varepsilon}|u(r,s)|^2\,ds\,dr
\lesssim
\|u\|_{H^1(\R^2_r;L^2(\R^2_s))}^2.
\]
Since \(|r|^{-\varepsilon}\le1\) on \(\R^2\setminus B_1(\mathbb{R}^2)\), we also control
\[
\int_{|r|>1}\int_{\R^2}|r|^{-\varepsilon}|u(r,s)|^2\,ds\,dr
\le
\|u\|_{L^2(\R^4)}^2.
\]
Therefore
\[
\int_{\R^2\times\R^2}|r|^{-\varepsilon}|u(r,s)|^2\,dr\,ds
\lesssim
\Bigl(\|u\|_{H^1(\R^2_r;L^2(\R^2_s))}^2+\|u\|_{L^2(\R^4)}^2\Bigr).
\]
By the Sobolev embedding \(H^1(\R^4)\hookrightarrow L^4(\R^4)\),
\[
\|u\|_{L^4(\R^4)}^2
\lesssim\|u\|_{H^1(\R^4)}^2
\lesssim
\langle u,(h_1+h_2+1)u\rangle.
\]
Finally, since \(0<\varepsilon<2\), there exist a constants $C_\varepsilon,C_\varepsilon'>0$ such that
\[
|x-y|^\varepsilon\le C_{\varepsilon}\bigl(1+|x|^\varepsilon+|y|^\varepsilon\bigr)
\le C'_{\varepsilon}\bigl(1+|x|^2+|y|^2\bigr),
\]
and therefore
\[
\int_{\R^2\times\R^2}|x-y|^\varepsilon |u(x,y)|^2\,dx\,dy
\le
C_\varepsilon'\langle u,(h_1+h_2+1)u\rangle.
\]
Putting the three estimates together yields
\[
\int_{\R^2\times\R^2}|\log|x-y||\,|u(x,y)|^2\,dx\,dy
\lesssim
\langle u,(h_1+h_2+1)u\rangle.
\]
Summing over all combinations of $x=x_i,y=x_j$ yields then 
\[
|W_n[\psi]|\lesssim B_n\langle\psi,(H_{0,n}+1)\psi\rangle.
\]
This ends the proof of \eqref{eq111}. For \eqref{eq222}, we may proceed for both $d=2,3$ as follows:
instead of bounding the interaction form entirely by the energy right away, we first apply the Cauchy--Schwarz inequality to the interaction potential. For $d=3$, using Hardy's inequality $\| |x_i-x_j|^{-1}\phi \| \lesssim \langle \phi, (h_i+h_j)\phi \rangle^{1/2}$, and for $d=2$ using the analogous interpolation from the Gagliardo--Nirenberg inequality derived above, we obtain the fractional form bound for any $\phi \in Q(H_{0,n})$:
\[
|W_n[\phi]| \lesssim B_n \|\phi\| \langle \phi, (H_{0,n}+1)\phi \rangle^{1/2}.
\]
We evaluate this bound on the high-energy projection $\phi = \Pi_M^\perp\psi$. Since $(H_{0,n}+1)\Pi_M^\perp\ge(\lambda_{M+1}+1)\Pi_M^\perp$, the norm decays as:
\[ 
\|\Pi_M^\perp\psi\| \le (\lambda_{M+1}+1)^{-1/2}\langle\psi,(H_{0,n}+1)\psi\rangle^{1/2}.
\]
Furthermore, since $[\Pi_M^\perp,H_{0,n}]=0$ and $\Pi_M^\perp \le 1$, we have $\langle \Pi_M^\perp\psi, (H_{0,n}+1)\Pi_M^\perp\psi\rangle \le \langle \psi, (H_{0,n}+1)\psi\rangle$. Substituting these into our fractional form bound yields
\begin{align}
|W_n[\Pi_M^\perp\psi]| &\lesssim B_n \|\Pi_M^\perp\psi\| \langle \Pi_M^\perp\psi, (H_{0,n}+1)\Pi_M^\perp\psi \rangle^{1/2}\nonumber\\
&\le B_n \left( (\lambda_{M+1}+1)^{-1/2} \langle \psi, (H_{0,n}+1)\psi \rangle^{1/2} \right) \langle \psi, (H_{0,n}+1)\psi \rangle^{1/2}\nonumber\\
&= (\lambda_{M+1}+1)^{-1/2}B_n\langle \psi,(H_{0,n}+1)\psi\rangle,\label{eq2222}
\end{align}
which ends the proof of \eqref{eq222}.


For \eqref{eq333}, we associate to \(W_n\) the sesquilinear form
\[
W_n(\varphi,\chi)
:=
\sum_{1\le i<j\le n}
\alpha_{n,i,j}\int_{(\R^d)^n}w_d(x_i-x_j)\,\varphi(x)\,\overline{\chi(x)}\,dx,
\]
so that \(W_n[\psi]=W_n(\psi,\psi)\). Hence, writing \(P:=\Pi_M\) and \(Q:=\Pi_M^\perp\),
\[
W_n[\psi]-W_{n,M}[\psi]
=
W_n[P\psi+Q\psi]-W_n[P\psi]
=
W_n[Q\psi]+2\Re W_n(P\psi,Q\psi).
\]
Similarly, if \(\widetilde W_n\) denotes the positive sesquilinear form with kernel
\(|w_d|\) and coefficients \(|\alpha_{n,i,j}|\), then combining the Cauchy-Schwarz inequality, \eqref{eq111} applied to \(P\psi\), and \eqref{eq222} applied to \(Q\psi\), we get
\begin{align}
|W_n(P\psi,Q\psi)|&\le \widetilde{W}_n[P\psi]^{1/2}\widetilde{W}_n[Q\psi]^{1/2}\nonumber\\
&\lesssim B_n(\lambda_{M+1}+1)^{-1/4}\langle \psi,(H_{0,n}+1)\psi\rangle.\nonumber
\end{align}
Finally, combining the last bound with the estimate for \(W_n[Q\psi]\) derived in \eqref{eq222} yields the claimed bound \eqref{eq333}.
\end{proof}

\end{document}
